% Part: turing-machines
% Chapter: machines-computations
% Section: turing-machines

\documentclass[../../../include/open-logic-section]{subfiles}

\begin{document}

\olfileid{tur}{mac}{tur}
\olsection{ट्यूरिंग यंत्रे}

\begin{explain}
ट्यूरिंग यंत्राची औपचारिक व्याख्या अमूर्त वाटते, पण प्रत्यक्षात
ती सोपी आहे: यंत्राचे कार्य कसे चालते हे निर्दिष्ट करण्यासाठी लागणारी
सर्व माहिती ती एका गणिती रचनेत एकत्र करते. यात (1) यंत्र कोणत्या
अवस्थांत असू शकते, (2) फितीवर कोणती चिन्हे असू शकतात, (3) यंत्राने
कोणत्या अवस्थेत सुरू व्हावे आणि (4) त्याचा सूचनासंच कोणता आहे,
या गोष्टी येतात.
\end{explain}

\begin{defn}[ट्यूरिंग यंत्र]
\emph{ट्यूरिंग यंत्र} $M$ म्हणजे पुढील घटक असलेली
$\langle Q, \Sigma, q_0,
\delta\rangle$ ही क्रमित चौकडी:
\begin{enumerate}
\item \emph{अवस्थांचा} सांत संच~$Q$,
\item $\TMendtape$ आणि $\TMblank$ यांचा समावेश असलेली सांत
  \emph{वर्णमाला} $\Sigma$,
\item \emph{प्रारंभिक अवस्था}~$q_0 \in Q$,
\item सांत \emph{सूचनासंच}~$\delta\colon Q \times \Sigma
  \pto Q \times \Sigma \times \{\TMleft, \TMright, \TMstay\}$.
\end{enumerate}
या आंशिक फलनाला~$\delta$ हे $M$ चे \emph{संक्रमण फलन} असेही म्हणतात.
\end{defn}

\begin{explain}
फीत फक्त एकाच दिशेने अनंत आहे, असे आपण मानतो. म्हणून फितीचे
डावे टोक दर्शवण्यासाठी $\TMendtape$ हे खास चिन्ह ठेवणे उपयुक्त ठरते.
त्यामुळे ट्यूरिंग यंत्राच्या कार्यक्रमाला फितीच्या बाहेर जाण्याचा
``धोका'' कधी आहे हे ओळखणे सोपे जाते. हे चिन्ह कधीही बदलून लिहिले
जात नाही, असे आपण गृहीत धरू शकलो असतो; म्हणजे
$\delta(q,\TMendtape)$ परिभाषित असेल, तर
$\delta(q,\TMendtape) = \tuple{q', \TMendtape, x}$ असेल.
काही पाठ्यपुस्तके असे करतात; आपण मात्र हे गृहीत धरत नाही. ट्यूरिंग
यंत्र बनवताना ते $\TMendtape$ कधीही बदलून लिहिणार नाही याची
सावधगिरी तुम्ही घेऊ शकता. शिवाय काही प्रसंगी असे बदलून लिहिण्यास
परवानगी दिल्यास सोयीची लवचीकता मिळते.
\end{explain}

\begin{ex}
\emph{सम यंत्र:} सम यंत्र औपचारिकरीत्या
$\tuple{Q, \Sigma, q_0, \delta}$ ही क्रमित चौकडी आहे. येथे
\begin{align*}
Q & = \{ q_0, q_1 \} \\
\Sigma & = \{ \TMendtape, \TMblank, \TMstroke \}, \\
\delta(q_0, \TMstroke) & = \tuple{q_1, \TMstroke, \TMright},\\
\delta(q_1, \TMstroke) & = \tuple{q_0, \TMstroke, \TMright},\\
\delta(q_1, \TMblank)  & = \tuple{q_1, \TMblank, \TMright}.
\end{align*}
\end{ex}

\end{document}
